\documentclass[10pt,a4paper,twocolumn]{article}
% --- 极简中文支持（pdfLaTeX 直接兼容，无需切换编译器）---
\usepackage{CJKutf8} % pdfLaTeX 下最可靠的中文支持

% --- 基础宏包（清理冗余，确保无冲突）---
\usepackage{amsmath, amsfonts, amssymb} % 数学公式
\usepackage{graphicx} % 插图
\usepackage{geometry} % 页面布局
\usepackage{listings} % 代码块
\usepackage{xcolor} % 颜色
\usepackage{hyperref} % 超链接
\usepackage{booktabs} % 表格
\usepackage{caption}
\usepackage{float}
\usepackage{subcaption} % 子图支持
\usepackage{multirow} % 表格多行
\usepackage{algorithm} % 算法环境
\usepackage{algorithmic} % 算法伪代码

% 数学符号宏定义
\newcommand{\fro}[1]{\left\|#1\right\|_F}
\newcommand{\tr}{\mathrm{tr}}
\newcommand{\argmin}{\mathop{\mathrm{argmin}}}
\newcommand{\R}{\mathbb{R}}

\title{网格参数化的局部-全局方法}
\author{
刘利刚$^1$ \quad 张磊$^1$ \quad 徐寅$^1$ \quad Craig Gotsman$^2$ \quad Steven J. Gortler$^3$ \\
$^1$ 浙江大学，中国 \quad $^2$ 以色列理工学院 \quad $^3$ 哈佛大学，美国 \\
\textit{ligangliu@zju.edu.cn,\ gotsman@cs.technion.ac.il,\ sjg@harvard.edu}
}
\date{}

\begin{document}

\begin{CJK*}{UTF8}{gbsn}
\maketitle

\begin{abstract}
本文提出一种保形的圆盘拓扑网格平面参数化方法，核心为**局部-全局算法**：将每个3D三角面片通过受限变换映射至平面（局部阶段），再通过稀疏线性系统拼接所有面片保证拓扑一致性（全局阶段）。
局部变换支持**相似变换（ASAP）**与**旋转变换（ARAP）**：ASAP等价于最小二乘共形映射（LSCM）；ARAP实现近等距参数化，保形性最优。
本文还提出单参数混合模型，实现保角与保面积的连续权衡；算法支持并行加速与矩阵预分解，效率优异。实验证明，该方法优于同期主流参数化算法。
\end{abstract}

\section{引言}
3D网格曲面参数化是计算机图形学核心技术，广泛应用于纹理映射、模型压缩、形状融合等任务。非可展曲面平面展开必然产生度量畸变，参数化的核心目标是**最小化角度、面积、长度畸变**，同时保证映射双射无重叠。
受网格ARAP变形算法启发，本文将参数化建模为**局部面片畸变最小化+全局网格一致性**的优化问题，首次将刚性变形思想迁移至网格参数化。

\section{相关工作}
主流参数化方法分为三类：
1. \textbf{线性重心坐标法}：基于Tutte定理，余切权重实现简单，但固定边界易产生畸变；自由边界LSCM实现共形映射，但存在三角面片翻转问题；
2. \textbf{非线性畸变优化}：MIPS、ABF++最小化角度/拉伸畸变，精度高但计算量大；
3. \textbf{曲率驱动共形映射}：将曲率集中于锥奇点，仅保角但面积控制能力弱。

\section{核心贡献}
1. 证明**ASAP参数化与LSCM算法数学等价**；
2. 提出ARAP参数化的局部-全局迭代算法，近等距保形，求解高效；
3. 设计单参数混合能量模型，连续权衡保角（ASAP）与保面积（ARAP）；
4. 局部阶段可GPU并行，全局阶段稀疏矩阵一次预分解复用，极致提速。

\section{局部-全局通用方法}
\subsection{能量函数定义}
设3D网格共$T$个三角面片，第$t$个面片面积为$A_t$，3D顶点坐标$\boldsymbol{x}_t$，2D参数域坐标$\boldsymbol{u}_t$，面片雅克比矩阵$J_t(\boldsymbol{u})$，局部变换矩阵$L_t$。
定义**弗罗贝尼乌斯范数畸变能量**：
\begin{equation}
E(\boldsymbol{u}, L)=\sum_{t=1}^{T} A_t \fro{J_t(\boldsymbol{u}) - L_t}
\end{equation}
展开为网格顶点坐标形式（余切权重）：
\begin{equation}
E = \frac{1}{2}\sum_{t=1}^T\sum_{i=0}^2 \cot\theta_t^i \left\| (\boldsymbol{u}_t^i-\boldsymbol{u}_t^{i+1}) - L_t(\boldsymbol{x}_t^i-\boldsymbol{x}_t^{i+1}) \right\|^2
\end{equation}
优化目标：
\begin{equation}
(\boldsymbol{u}, L) = \argmin_{\boldsymbol{u},L} E(\boldsymbol{u},L) \quad \text{s.t.} \ L_t \in M
\end{equation}
其中$M$为允许的变换集合，$L_t$为辅助变量，最终仅需求解2D坐标$\boldsymbol{u}$。

\subsection{最优变换矩阵求解（Procrustes分析）}
对雅克比矩阵做**带符号奇异值分解（SVD）**：
$$ J = U\Sigma V^T, \quad \Sigma=\begin{pmatrix}\sigma_1&0\\0&\sigma_2\end{pmatrix} $$
约束$\det(UV^T)>0$避免方向反转，最优变换：
1. 旋转矩阵（ARAP）：$L=UV^T$（奇异值置1）；
2. 相似矩阵（ASAP）：$L=\frac{\sigma_1+\sigma_2}{2}UV^T$（奇异值取均值）。

\subsection{ASAP（尽可能相似）参数化}
限定变换为**相似矩阵**（保角）：
$$ M=\left\{ \begin{pmatrix}a&b\\-b&a\end{pmatrix} \mid a,b\in\R \right\} $$
最小化共形能量，**严格等价于LSCM**：
\begin{equation}
E(\boldsymbol{u}) = \sum_{t=1}^T A_t (\sigma_{1,t}-\sigma_{2,t})^2
\end{equation}
固定网格最远两点避免平凡解（全顶点坍缩）。

\subsection{ARAP（尽可能刚性）参数化}
限定变换为**纯旋转矩阵**（保面积/长度）：
$$ M=\left\{ \begin{pmatrix}\cos\theta&\sin\theta\\-\sin\theta&\cos\theta\end{pmatrix} \mid \theta\in[0,2\pi) \right\} $$
最小化等距畸变能量：
\begin{equation}
E(\boldsymbol{u}) = \sum_{t=1}^T A_t \left[ (\sigma_{1,t}-1)^2 + (\sigma_{2,t}-1)^2 \right]
\end{equation}
该能量非线性，需迭代求解。

\subsection{局部-全局迭代算法}
\subsubsection{局部阶段}
固定2D坐标$\boldsymbol{u}$，对每个面片做SVD分解，求解最优旋转矩阵$L_t$，**面片间独立可并行**。
\subsubsection{全局阶段}
固定$L_t$，能量对$\boldsymbol{u}$二次凸，梯度置零构建**稀疏线性系统**：
$$ \sum_{(i,j)\in he} \cot\theta_{ij} \left[ (\boldsymbol{u}_i-\boldsymbol{u}_j) - L_{t(i,j)}(\boldsymbol{x}_i-\boldsymbol{x}_j) \right] = 0 $$
矩阵仅依赖3D网格拓扑，**一次Cholesky预分解，全迭代复用**。
\subsubsection{后处理}
采用凸虚拟边界算法修复少量三角面片翻转问题。

\subsection{初始化策略}
单边界网格用Floater保形参数化，多边界网格用LSCM初始化；算法对初始值不敏感，快速稳定收敛。

\section{混合参数化模型}
引入权重$\lambda\in[0,+\infty)$，统一ASAP与ARAP，能量函数：
\begin{equation}
E = \frac{1}{2}\sum_{t=1}^T\left[ \sum_{i=0}^2 \cot\theta_t^i \fro{\nabla e_t^i} + \lambda(a_t^2+b_t^2-1)^2 \right]
\end{equation}
其中：
$$ \nabla e_t^i = (\boldsymbol{u}_t^i-\boldsymbol{u}_t^{i+1}) - \begin{pmatrix}a_t&b_t\\-b_t&a_t\end{pmatrix}(\boldsymbol{x}_t^i-\boldsymbol{x}_t^{i+1}) $$
- $\lambda=0$：纯ASAP（保角优先）；
- $\lambda\to\infty$：纯ARAP（保面积优先）。

局部阶段解析求解三次方程，全局阶段复用稀疏线性系统，效率与ARAP持平。

\section{实验结果与对比}
\subsection{畸变量化指标}
基于雅克比奇异值定义角度/面积畸变：
\begin{equation}
D^\text{angle} = \sum_t \rho_t \left( \frac{\sigma_{t1}}{\sigma_{t2}} + \frac{\sigma_{t2}}{\sigma_{t1}} \right)
\end{equation}
\begin{equation}
D^\text{area} = \sum_t \rho_t \left( \sigma_{t1}\sigma_{t2} + \frac{1}{\sigma_{t1}\sigma_{t2}} \right)
\end{equation}
权重$\rho_t = A_t / \sum A_t$，值越小畸变越低。

\subsection{实验对比}
对比LSCM、ABF++、LABF、IC、CP算法，**ARAP在面积畸变上全面最优**，角度畸变仅微小增加；算法10次迭代内收敛，支持多边界网格，效率与ABF++相当。

\begin{table*}[htbp]
\centering
\caption{不同参数化方法畸变对比}
\resizebox{\textwidth}{!}{
\begin{tabular}{lcccccc|cccccc}
\toprule
\multirow{2}{*}{模型} & \multicolumn{6}{c|}{角度畸变 $D^\text{angle}$} & \multicolumn{6}{c}{面积畸变 $D^\text{area}$} \\
& LSCM & dABF & LABF & IC & CP & ARAP & LSCM & dABF & LABF & IC & CP & ARAP \\
\midrule
石像鬼 & 2.00 & 2.00 & 2.00 & 2.05 & 2.00 & 2.06 & 88.14 & 2.64 & 2.64 & 2.67 & 2.64 & 2.05 \\
伊西斯 & 2.05 & 2.08 & 2.12 & 3.09 & 2.29 & 2.19 & 15.36 & 8.37 & 9.12 & 3.91 & 11.94 & 2.11 \\
球体 & 2.02 & 2.06 & 2.05 & 2.06 & 2.36 & 2.07 & 2.40 & 2.35 & 2.34 & 2.42 & 3.18 & 2.04 \\
奶牛 & 2.01 & 2.01 & 2.01 & 2.46 & 2.01 & 2.03 & 30.09 & 2.19 & 2.22 & 2.51 & 2.18 & 2.03 \\
甲虫 & 2.00 & 2.00 & — & 2.02 & 2.00 & 2.01 & 2.08 & 2.09 & — & 2.14 & 2.10 & 2.01 \\
\bottomrule
\end{tabular}
}
\end{table*}

\section{讨论与结论}
本文提出的局部-全局框架统一了共形与刚性网格参数化，兼具**高精度、高效率、高灵活性**：局部最小化面片畸变，全局保证拓扑一致性；支持并行计算与矩阵预分解，适配工业级应用。
该方法仅适用于圆盘拓扑网格，未来可扩展至高亏格网格与带约束参数化场景。

\section{致谢}
感谢Yongliang Yang、Miri Ben-Chen提供算法实验结果；感谢Hugues Hoppe的技术指导。本工作受国家自然科学基金、美以双边科学基金资助。

\section{附录A 能量等价性证明}
最小化能量 $E(\boldsymbol{u},L)=\sum A_t \fro{J_t-L_t}$，固定$\boldsymbol{u}$时最优$L_t$由SVD得：
$$ L_t=U_t K_t V_t^T, \quad K_t=\mathrm{diag}(k_{1t},k_{2t}) $$
代入能量展开：
$$ E(\boldsymbol{u}) = \sum_{t=1}^T A_t \left[ (\sigma_{1t}-k_{1t})^2 + (\sigma_{2t}-k_{2t})^2 \right] $$
1. ARAP：$k_{1t}=k_{2t}=1$，得等距能量；
2. ASAP：$k_{1t}=k_{2t}=(\sigma_{1t}+\sigma_{2t})/2$，得LSCM共形能量。

\section{附录B 混合模型解析解}
对能量求偏导置零，得方程组：
\begin{equation}
C_1 a + 2\lambda a(a^2+b^2-1) = C_2
\end{equation}
\begin{equation}
C_1 b + 2\lambda b(a^2+b^2-1) = C_3
\end{equation}
化简得单变量三次方程：
\begin{equation}
\frac{2\lambda(C_2^2+C_3^2)}{C_2^2}a^3 + (C_1-2\lambda)a - C_2 = 0
\end{equation}
特例解：
- $\lambda=0$：$a=C_2/C_1,\ b=C_3/C_1$；
- $\lambda\to\infty$：$a=\pm C_2/\sqrt{C_2^2+C_3^2},\ b=\pm C_3/\sqrt{C_2^2+C_3^2}$。

\begin{thebibliography}{99}
\bibitem{SA07} Sorkine O, Alexa M. 尽可能刚性曲面建模. 欧洲图形学几何处理研讨会, 2007.
\bibitem{LPRM02} Lévy B 等. 用于纹理图集生成的最小二乘共形映射. ACM SIGGRAPH, 2002.
\bibitem{SLMB05} Sheffer A 等. ABF++：快速鲁棒的基于角度的扁平化算法. ACM图形学汇刊, 2005.
\bibitem{YKL08} Yang Y 等. 基于逆曲率图的最优曲面参数化. IEEE可视化与计算机图形学汇刊, 2008.
\end{thebibliography}

\end{document}